\subsection{Improper Integrals}

\begin{definition} \label{def:improper-riemann-integral} \hlt{Improper Riemann Integrals}\\
Let $f\in\mathcal{R}$ on every compact subinterval of its domain.\\
Set $\int_a^b f(x)\,dx\defeq\lim\limits_{d\to b^-}\int_a^d f(x)\,dx$ at a right endpoint, and $\int_a^b f(x)\,dx\defeq\lim\limits_{c\to a^+}\int_c^b f(x)\,dx$ at a left endpoint.\\
If both endpoints are improper, choose $\xi\in(a,b)$ and set $\int_a^b f(x)\,dx \defeq \int_a^\xi f(x)\,dx+\int_\xi^b f(x)\,dx$, provided both one-sided integrals exist and the sum is not indeterminate.\\
A finite value means the improper integral \hlt{converges}; otherwise it diverges to $+\infty$ or $-\infty$.
\end{definition}

\begin{figure}[H]
  \centering
  \begin{tikzpicture}[
  x=0.78cm,
  y=0.72cm,
  >=Stealth,
  every node/.style={font=\scriptsize},
  axis/.style={->, black!55, line width=0.45pt},
  curve/.style={black!85, line width=0.9pt},
  area/.style={fill=gray!16, draw=none},
  guide/.style={black!30, dashed, line width=0.4pt}
]
  \begin{scope}
    \draw[axis] (-0.2,0) -- (5.05,0) node[right] {$x$};
    \draw[axis] (0,-0.15) -- (0,3.15);
    \path[area]
      (0,0) -- plot[domain=0:3.65, samples=100] (\x,{0.42+0.28*\x+0.07/((4.28-\x)*(4.28-\x))}) -- (3.65,0) -- cycle;
    \draw[curve, domain=0:4.02, samples=130]
      plot (\x,{0.42+0.28*\x+0.07/((4.28-\x)*(4.28-\x))});
    \draw[guide] (3.65,0) -- (3.65,2.35);
    \draw[guide] (4.28,0) -- (4.28,3.00);
    \node[below] at (0,0) {$a$};
    \node[below] at (3.65,0) {$d$};
    \node[below] at (4.28,0) {$b$};
    \node at (2.25,-0.55) {$d\to b^-$};
  \end{scope}

  \begin{scope}[xshift=6.4cm]
    \draw[axis] (-0.2,0) -- (5.7,0) node[right] {$x$};
    \draw[axis] (0,-0.15) -- (0,3.15);
    \path[area]
      (0,0) -- plot[domain=0:4.35, samples=120] (\x,{2.35/(1+\x)+0.18}) -- (4.35,0) -- cycle;
    \draw[curve, domain=0:5.25, samples=150]
      plot (\x,{2.35/(1+\x)+0.18});
    \draw[guide] (4.35,0) -- (4.35,0.62);
    \node[below] at (0,0) {$a$};
    \node[below] at (4.35,0) {$d$};
    \node[below] at (5.45,0) {$\infty$};
    \node at (2.80,-0.55) {$d\to\infty$};
  \end{scope}
\end{tikzpicture}

  \caption{Improper integrals are limits of ordinary integrals over compact truncations.}
  \label{fig:improper-integral-truncations}
\end{figure}

\begin{theorem} \label{thm:nonnegative-improper-integral-supremum} \hlt{Nonnegative Improper Integrals}\\
Let $f\geq0$ and $f\in\mathcal{R}$ on every $[a,d]$, $a<d<b$. Then $\int_a^b f(x)\,dx=\sup\left\{\int_a^d f(x)\,dx:a<d<b\right\}$.\\
Thus $\int_a^b f$ exists in $[0,+\infty]$ and converges exactly when this supremum is finite.
\begin{proof}
For $a<d_1<d_2<b$, $\int_a^{d_2}f-\int_a^{d_1}f=\int_{d_1}^{d_2}f\geq0$.\\
Hence $d\mapsto\int_a^d f$ is increasing, so its extended limit exists and equals its supremum.\\
\hlt{Structural Commentary}: Nonnegative area can only accumulate.\\
\hlt{Key Insight}: Improper convergence for $f\geq0$ is boundedness of all finite truncations.
\end{proof}
\end{theorem}

\begin{perspective} \label{per:cauchy-principal-value-not-improper-integral} \hlt{Principal Value Is Different}\\
Improper integrals over $(-\infty,\infty)$ require separate one-sided limits.\\
A symmetric limit such as $\lim_{A\to\infty}\int_{-A}^{A}\sin x\,dx$ may still exist (Cauchy principal value, not Definition \ref{def:improper-riemann-integral}).
\end{perspective}

\begin{definition} \label{def:improper-riemann-stieltjes-real-line} \hlt{Improper Riemann--Stieltjes Integrals on $\R$}\\
Let $\alpha:\R\to\R$ be increasing and $f\in\mathcal{R}(\alpha)$ on every finite interval.\\
Define, when the limits exist, $\int_0^\infty f\,d\alpha\defeq\lim_{b\to\infty}\int_0^b f\,d\alpha$; and $\int_{-\infty}^{0} f\,d\alpha\defeq\lim_{a\to-\infty}\int_a^0 f\,d\alpha$.\\
If their sum is not indeterminate, set $\int_{-\infty}^{\infty}f\,d\alpha \defeq \int_{-\infty}^{0}f\,d\alpha+\int_0^\infty f\,d\alpha$.\\
Finite value means $f$ is \hlt{$\alpha$-integrable on $\R$}. Write $\alpha(\pm\infty)\defeq\lim_{t\to\pm\infty}\alpha(t)$ when these extended limits are used.
\end{definition}

\begin{theorem} \label{thm:improper-riemann-stieltjes-nonnegative-bounded-mass} \hlt{Nonnegative and Bounded-Mass Stieltjes Integrals}\\
Let $\alpha:\R\to\R$ be increasing, and suppose $f\in\mathcal{R}(\alpha)$ on every finite interval.
\begin{enumerate}[label=(\roman*)]
\setlength{\itemsep}{0pt}
\item If $f\geq0$, then either $f$ is $\alpha$-integrable on $\R$ or $\int_{-\infty}^{\infty}f\,d\alpha=+\infty$.
\item If $-\infty<\alpha(-\infty)<\alpha(\infty)<+\infty$ and $f$ is bounded, then $f$ is $\alpha$-integrable on $\R$.
\end{enumerate}
\begin{proof}
\begin{enumerate}[label=(\roman*)]
\setlength{\itemsep}{0pt}
\item
By additivity and order preservation, $\int_0^b f\,d\alpha$ and $\int_a^0 f\,d\alpha$ are monotone under outward truncation.\\
Hence both extended one-sided limits exist, so their sum is finite or $+\infty$.
\item
Let $\abs{f}\leq M$. For finite $a<b$, Theorem \ref{thm:riemann-stieltjes-integral-properties} gives $\abs{\int_a^b f\,d\alpha}\leq M\bigl[\alpha(b)-\alpha(a)\bigr]$.\\
Since $\alpha(\infty)$ and $\alpha(-\infty)$ are finite, the right and left tails have vanishing $\alpha$-increment.\\
Thus the one-sided truncations are Cauchy and have finite limits.
\end{enumerate}
\hlt{Structural Commentary}: Positivity gives monotone limits; finite mass bounds tails.\\
\hlt{Key Insight}: A bounded increasing integrator behaves like finite total mass.
\end{proof}
\end{theorem}
